\section{Related Work and Background}
\label{sec:background}

This section reviews TinyML constraints and recent Mamba literature relevant to MCU deployment. We summarize common TinyML optimizations—focusing on quantization, pruning, and distillation—and detail prior deployment attempts that motivate our contributions.

\subsection{TinyML Constraints and Optimization Strategies}

Given the limited computational throughput of MCUs, minimizing latency is a critical operational priority for real-time edge processing (e.g., audio or accelerometer streams). Many complex deep learning architectures rely on operations such as exponentiation (e.g., Softmax, SiLU) or trigonometric functions (e.g., Tanh). On microcontrollers, computing these functions non-natively introduces substantial execution overhead. Replacing them with hardware-friendly alternatives like ReLU can provide up to a 25$\times$ computational speedup under resource-constrained conditions \cite{cassimonDesigningResourceconstrainedNeural2020}.

The lack of robust or fast hardware floating-point units (FPUs) on many MCUs necessitates post-training or aware quantization. Quantization reduces the precision of weights and activations from 32-bit floating-point (FP32) formats to lower-bit fixed-point representations, such as 8-bit or 4-bit integers (INT8/INT4). While this strategy substantially reduces the model's memory footprint, these gains often incur an accuracy penalty. Furthermore, the final memory savings are not always directly proportional to the reduced bit-width due to underlying software graph and format conversion overheads \cite{dokicInferenceSpeedQuantisation2020}. Careless precision reduction can yield severe accuracy drops, particularly for neural network layers that were not structurally designed with quantization bounds in mind \cite{xuMambaQuantQuantizingMamba2025}.

Other optimization mechanisms include network pruning, knowledge distillation, and low-rank factorization \cite{suwannaphongOptimisingTinyMLQuantization2025}. Pruning algorithms eliminate redundant weight channels, whereas knowledge distillation transfers latent representations from a larger teacher network into a compact student model. Although highly effective, applying these methods to alternative structures is often complex, and achieved compression ratios vary significantly depending on the target task and network layout \cite{shihabEfficientUnstructuredPruning}.

\subsection{Mamba Architecture and Edge Deployment Efforts}

Most prior literature exploring Mamba model compression has focused on mobile or embedded GPUs rather than MCU-class systems. For example, the Physics-Guided Tiny-Mamba Transformer \cite{liPhysicsGuidedTinyMambaTransformer2026} compresses its topology down to $\sim$0.5M parameters for deployment on the NVIDIA Jetson Nano; Tiny-ViM \cite{maTinyViMFrequencyDecoupling2025} targets vision tasks at a scale of $\sim$5M parameters; and FEMBA \cite{tegonFEMBAEdgePhysiologicallyAware2026a} utilizes aggressive quantization targeting a $\sim$7.8M parameter model. While these works demonstrate that standard optimizations like quantization, pruning, and distillation are viable for State-Space Models, the resulting architectures remain orders of magnitude too large for typical microcontroller memory budgets.

Efforts such as LightMamba \cite{wangLightMambaResourceEfficientDeepLearning2026} explore Selective State-Space Models for specialized, resource-constrained tasks like ultra-wideband positioning, but they stop short of evaluating physical microcontroller deployments. MambaLite-Micro \cite{xuMambaLiteMicroMemoryOptimizedMamba2025} is the only previous work that we are aware of that actually achieves successful execution on a microcontroller. It uses ESP32-S3 and STM32 MCUs and a custom C mamba implementation, but its strict adherence to PyTorch-style kernels complicates deep architectural modifications. Consequently, deployment tooling and structural co-design methodologies for Mamba models on MCUs remain immature.

\subsection{Research Gap}

Two primary gaps motivate this work. First, while many TinyML optimization frameworks exist, their algorithmic behavior and stability when applied to Mamba's input-dependent recurrence are not well understood on MCU-class hardware. Second, while initial C-based engines demonstrate basic execution feasibility, they do not explore explicit graph modifications or hardware-aware activation substitutions to mitigate MCU computational bottlenecks. This paper bridges these gaps by profiling physical hardware deployments, isolating structural bottlenecks, and evaluating targeted architectural changes to make Mamba models truly viable on microcontrollers.
